%% PATRÓN · Selector: cada columna es un eje de decisión, cada chip una opción. Se enciende una
\providecolor{textochip}{named}{white} % por si la charla tiene un estilo/ anterior a 1.1.0
%% opción por columna y un trazo las une: la «receta» de un caso. Con todo vacío, solo las opciones
%% (la slide que explica los ejes una vez). Con \only<n>{...} se recorren varios casos.
%% Uso: %% PATRÓN · Selector: cada columna es un eje de decisión, cada chip una opción. Se enciende una
\providecolor{textochip}{named}{white} % por si la charla tiene un estilo/ anterior a 1.1.0
%% opción por columna y un trazo las une: la «receta» de un caso. Con todo vacío, solo las opciones
%% (la slide que explica los ejes una vez). Con \only<n>{...} se recorren varios casos.
%% Uso: %% PATRÓN · Selector: cada columna es un eje de decisión, cada chip una opción. Se enciende una
\providecolor{textochip}{named}{white} % por si la charla tiene un estilo/ anterior a 1.1.0
%% opción por columna y un trazo las une: la «receta» de un caso. Con todo vacío, solo las opciones
%% (la slide que explica los ejes una vez). Con \only<n>{...} se recorren varios casos.
%% Uso: %% PATRÓN · Selector: cada columna es un eje de decisión, cada chip una opción. Se enciende una
\providecolor{textochip}{named}{white} % por si la charla tiene un estilo/ anterior a 1.1.0
%% opción por columna y un trazo las une: la «receta» de un caso. Con todo vacío, solo las opciones
%% (la slide que explica los ejes una vez). Con \only<n>{...} se recorren varios casos.
%% Uso: \input{diagramas/selector.tex} en el preámbulo, y en el frame \selector{imp}{asgs}{din}{}{nombre}.
%% Original con 6 columnas: charla01/ingredientes.tex. Para cambiar columnas, editar \selector abajo.
\makeatletter
\newcommand\sel@opcion[4]{% #1 color, #2 clave, #3 seleccionada, #4 etiqueta
  \ifnum\pdfstrcmp{#2}{#3}=0
    \node[fill=#1,text=textochip,font=\footnotesize\bfseries\boldmath,sel chip] (n) at (\sel@x,\sel@y) {#4};
    \ifx\sel@ultimo\empty\else
      \begin{pgfonlayer}{background}\draw[tinta!45!fondo,line width=1.1pt] (\sel@ultimo) -- (n);\end{pgfonlayer}\fi
    \edef\sel@ultimo{\sel@x,\sel@y}%
  \else
    \node[fill=#1!11!fondo,text=#1!75!tinta,font=\footnotesize,sel chip] at (\sel@x,\sel@y) {#4};
  \fi
  \pgfmathsetmacro\sel@y{\sel@y-0.76}}
\newcommand\sel@columna[3]{% #1 color, #2 índice, #3 título
  \pgfmathsetmacro\sel@x{1.45+3.4*#2}\def\sel@y{-0.75}%
  \node[anchor=base,font=\small\bfseries,text=#1] at (\sel@x,0) {#3};}
\newcommand\selector[5]{%
  \begin{tikzpicture}[sel chip/.style={rounded corners=2pt,minimum width=2.9cm,minimum height=0.6cm,inner sep=1.5pt}]
    \def\sel@ultimo{}%
    \sel@columna{serieA}{0}{Esquema}
      \sel@opcion{serieA}{imp}{#1}{implícito}
      \sel@opcion{serieA}{exp}{#1}{explícito}
    \sel@columna{serieB}{1}{Estabilización}
      \sel@opcion{serieB}{no}{#2}{ninguna}
      \sel@opcion{serieB}{asgs}{#2}{ASGS}
      \sel@opcion{serieB}{oss}{#2}{OSS}
    \sel@columna{serieC}{2}{Subescalas}
      \sel@opcion{serieC}{qs}{#3}{cuasiestáticas}
      \sel@opcion{serieC}{din}{#3}{dinámicas}
    \sel@columna{serieD}{3}{Modelo LES}
      \sel@opcion{serieD}{no}{#4}{ninguno}
      \sel@opcion{serieD}{vreman}{#4}{Vreman}
    \node[anchor=west,inner sep=0pt] at (0.08,-3.3) {%
      \if\relax\detokenize{#5}\relax{\small\color{tenue}Un caso = una opción de cada columna}%
      \else{\color{tenue}\LARGE$=$}\enspace{\LARGE\caso{#5}}\fi};
  \end{tikzpicture}}
\makeatother
 en el preámbulo, y en el frame \selector{imp}{asgs}{din}{}{nombre}.
%% Original con 6 columnas: charla01/ingredientes.tex. Para cambiar columnas, editar \selector abajo.
\makeatletter
\newcommand\sel@opcion[4]{% #1 color, #2 clave, #3 seleccionada, #4 etiqueta
  \ifnum\pdfstrcmp{#2}{#3}=0
    \node[fill=#1,text=textochip,font=\footnotesize\bfseries\boldmath,sel chip] (n) at (\sel@x,\sel@y) {#4};
    \ifx\sel@ultimo\empty\else
      \begin{pgfonlayer}{background}\draw[tinta!45!fondo,line width=1.1pt] (\sel@ultimo) -- (n);\end{pgfonlayer}\fi
    \edef\sel@ultimo{\sel@x,\sel@y}%
  \else
    \node[fill=#1!11!fondo,text=#1!75!tinta,font=\footnotesize,sel chip] at (\sel@x,\sel@y) {#4};
  \fi
  \pgfmathsetmacro\sel@y{\sel@y-0.76}}
\newcommand\sel@columna[3]{% #1 color, #2 índice, #3 título
  \pgfmathsetmacro\sel@x{1.45+3.4*#2}\def\sel@y{-0.75}%
  \node[anchor=base,font=\small\bfseries,text=#1] at (\sel@x,0) {#3};}
\newcommand\selector[5]{%
  \begin{tikzpicture}[sel chip/.style={rounded corners=2pt,minimum width=2.9cm,minimum height=0.6cm,inner sep=1.5pt}]
    \def\sel@ultimo{}%
    \sel@columna{serieA}{0}{Esquema}
      \sel@opcion{serieA}{imp}{#1}{implícito}
      \sel@opcion{serieA}{exp}{#1}{explícito}
    \sel@columna{serieB}{1}{Estabilización}
      \sel@opcion{serieB}{no}{#2}{ninguna}
      \sel@opcion{serieB}{asgs}{#2}{ASGS}
      \sel@opcion{serieB}{oss}{#2}{OSS}
    \sel@columna{serieC}{2}{Subescalas}
      \sel@opcion{serieC}{qs}{#3}{cuasiestáticas}
      \sel@opcion{serieC}{din}{#3}{dinámicas}
    \sel@columna{serieD}{3}{Modelo LES}
      \sel@opcion{serieD}{no}{#4}{ninguno}
      \sel@opcion{serieD}{vreman}{#4}{Vreman}
    \node[anchor=west,inner sep=0pt] at (0.08,-3.3) {%
      \if\relax\detokenize{#5}\relax{\small\color{tenue}Un caso = una opción de cada columna}%
      \else{\color{tenue}\LARGE$=$}\enspace{\LARGE\caso{#5}}\fi};
  \end{tikzpicture}}
\makeatother
 en el preámbulo, y en el frame \selector{imp}{asgs}{din}{}{nombre}.
%% Original con 6 columnas: charla01/ingredientes.tex. Para cambiar columnas, editar \selector abajo.
\makeatletter
\newcommand\sel@opcion[4]{% #1 color, #2 clave, #3 seleccionada, #4 etiqueta
  \ifnum\pdfstrcmp{#2}{#3}=0
    \node[fill=#1,text=textochip,font=\footnotesize\bfseries\boldmath,sel chip] (n) at (\sel@x,\sel@y) {#4};
    \ifx\sel@ultimo\empty\else
      \begin{pgfonlayer}{background}\draw[tinta!45!fondo,line width=1.1pt] (\sel@ultimo) -- (n);\end{pgfonlayer}\fi
    \edef\sel@ultimo{\sel@x,\sel@y}%
  \else
    \node[fill=#1!11!fondo,text=#1!75!tinta,font=\footnotesize,sel chip] at (\sel@x,\sel@y) {#4};
  \fi
  \pgfmathsetmacro\sel@y{\sel@y-0.76}}
\newcommand\sel@columna[3]{% #1 color, #2 índice, #3 título
  \pgfmathsetmacro\sel@x{1.45+3.4*#2}\def\sel@y{-0.75}%
  \node[anchor=base,font=\small\bfseries,text=#1] at (\sel@x,0) {#3};}
\newcommand\selector[5]{%
  \begin{tikzpicture}[sel chip/.style={rounded corners=2pt,minimum width=2.9cm,minimum height=0.6cm,inner sep=1.5pt}]
    \def\sel@ultimo{}%
    \sel@columna{serieA}{0}{Esquema}
      \sel@opcion{serieA}{imp}{#1}{implícito}
      \sel@opcion{serieA}{exp}{#1}{explícito}
    \sel@columna{serieB}{1}{Estabilización}
      \sel@opcion{serieB}{no}{#2}{ninguna}
      \sel@opcion{serieB}{asgs}{#2}{ASGS}
      \sel@opcion{serieB}{oss}{#2}{OSS}
    \sel@columna{serieC}{2}{Subescalas}
      \sel@opcion{serieC}{qs}{#3}{cuasiestáticas}
      \sel@opcion{serieC}{din}{#3}{dinámicas}
    \sel@columna{serieD}{3}{Modelo LES}
      \sel@opcion{serieD}{no}{#4}{ninguno}
      \sel@opcion{serieD}{vreman}{#4}{Vreman}
    \node[anchor=west,inner sep=0pt] at (0.08,-3.3) {%
      \if\relax\detokenize{#5}\relax{\small\color{tenue}Un caso = una opción de cada columna}%
      \else{\color{tenue}\LARGE$=$}\enspace{\LARGE\caso{#5}}\fi};
  \end{tikzpicture}}
\makeatother
 en el preámbulo, y en el frame \selector{imp}{asgs}{din}{}{nombre}.
%% Original con 6 columnas: charla01/ingredientes.tex. Para cambiar columnas, editar \selector abajo.
\makeatletter
\newcommand\sel@opcion[4]{% #1 color, #2 clave, #3 seleccionada, #4 etiqueta
  \ifnum\pdfstrcmp{#2}{#3}=0
    \node[fill=#1,text=textochip,font=\footnotesize\bfseries\boldmath,sel chip] (n) at (\sel@x,\sel@y) {#4};
    \ifx\sel@ultimo\empty\else
      \begin{pgfonlayer}{background}\draw[tinta!45!fondo,line width=1.1pt] (\sel@ultimo) -- (n);\end{pgfonlayer}\fi
    \edef\sel@ultimo{\sel@x,\sel@y}%
  \else
    \node[fill=#1!11!fondo,text=#1!75!tinta,font=\footnotesize,sel chip] at (\sel@x,\sel@y) {#4};
  \fi
  \pgfmathsetmacro\sel@y{\sel@y-0.76}}
\newcommand\sel@columna[3]{% #1 color, #2 índice, #3 título
  \pgfmathsetmacro\sel@x{1.45+3.4*#2}\def\sel@y{-0.75}%
  \node[anchor=base,font=\small\bfseries,text=#1] at (\sel@x,0) {#3};}
\newcommand\selector[5]{%
  \begin{tikzpicture}[sel chip/.style={rounded corners=2pt,minimum width=2.9cm,minimum height=0.6cm,inner sep=1.5pt}]
    \def\sel@ultimo{}%
    \sel@columna{serieA}{0}{Esquema}
      \sel@opcion{serieA}{imp}{#1}{implícito}
      \sel@opcion{serieA}{exp}{#1}{explícito}
    \sel@columna{serieB}{1}{Estabilización}
      \sel@opcion{serieB}{no}{#2}{ninguna}
      \sel@opcion{serieB}{asgs}{#2}{ASGS}
      \sel@opcion{serieB}{oss}{#2}{OSS}
    \sel@columna{serieC}{2}{Subescalas}
      \sel@opcion{serieC}{qs}{#3}{cuasiestáticas}
      \sel@opcion{serieC}{din}{#3}{dinámicas}
    \sel@columna{serieD}{3}{Modelo LES}
      \sel@opcion{serieD}{no}{#4}{ninguno}
      \sel@opcion{serieD}{vreman}{#4}{Vreman}
    \node[anchor=west,inner sep=0pt] at (0.08,-3.3) {%
      \if\relax\detokenize{#5}\relax{\small\color{tenue}Un caso = una opción de cada columna}%
      \else{\color{tenue}\LARGE$=$}\enspace{\LARGE\caso{#5}}\fi};
  \end{tikzpicture}}
\makeatother
